\begin{statement}[C.7.7; Known with derivation: two unconditional normalized values]
Let $E/\overline{\Q}$ have an algebraic invariant differential,
and let $u$ be an elliptic logarithm of an algebraic nontorsion
point $P$. Then
\[
 \sigma(u)\exp\left(-\frac12u\zeta(u)\right)
\]
is transcendental. If $E$ has CM and $s_2$ is normalized as in
C.7.6, then $\zeta(u)-s_2u$ is also transcendental. The first
assertion concerns the displayed holomorphic normalization;
it does not assert transcendence of either $N(P)$ or
$\theta(u)=\sigma(u)e^{-s_2u^2/2}$.
\end{statement}
\begin{proof}
We use \citet[Theorem 1 and Corollary 1]{waldschmidt1979sigma},
whose complete proof was opened in the translation hosted by its
author. Its Theorem 1 states that, for algebraic invariants
$g_2,g_3$, algebraic elliptic points $\exp_Eu$, $\exp_Ev$, and
$\beta\in\overline{\Q}$, the number
\[
 \frac{\sigma(u+v)}{\sigma(u)\sigma(v)}
                      \exp((\beta-\zeta(v))u)
\]
is transcendental if $\exp_Ev$ is nontorsion and $u,u+v$
are not lattice points. Here ``algebraic elliptic point''
means that the Weierstrass coordinates are algebraic; it does
not mean that its logarithm is algebraic. Take $v=u$ and
$\beta=0$. All hypotheses hold because $P$ is nontorsion.
The sigma duplication formula identifies the displayed number
with
\[
 -\wp'(u)\sigma(u)^2e^{-u\zeta(u)}.
\]
The factor $-\wp'(u)$ is nonzero and algebraic. Therefore
$\sigma(u)e^{-u\zeta(u)/2}$ is transcendental, as also stated
explicitly after Corollary 1 of the cited paper.

For the second assertion we apply the analytic subgroup theorem
in \citet[Theorem 6.1]{huber2022}. In the algebraic universal
vector extension $\widetilde E$ a logarithm of an algebraic lift
of $P$ has coordinates $(u,\zeta(u)-b)$, where $b$ is an algebraic
fibre coordinate, with a possible algebraic linear change of the
second coordinate. If $b=\zeta(u)-s_2u$ were algebraic, this
logarithm would belong to the algebraic line
$\{(z,s_2z):z\in\C\}$. The theorem applies to the connected
commutative algebraic group $\widetilde E$ over
$\overline{\Q}$ and this line, defined over $\overline{\Q}$ by
C.7.6; its exponential contains a nonzero algebraic lift.
It would provide a connected one-dimensional algebraic subgroup
whose Lie algebra is that line. Its projection to $E$ is
nontrivial. A one-dimensional affine algebraic group has no
nonconstant algebraic-group homomorphism to an elliptic curve,
so the subgroup is an elliptic curve isogenous to $E$.
Its inclusion would split the universal vector extension after
that isogeny. This is impossible: the nonzero extension class in
$\operatorname{Ext}^1(E,\mathbb G_a)=H^1(E,\mathcal O_E)$
remains nonzero under isogeny, since pullback is invertible in
characteristic zero.
\end{proof}
